\section{Conflicting source meaning and calibrated candidate histories}
\label{s6:linked-histories}
Consider a deposit feed with business date, warehouse arrival date and
revision number, plus a separate account/customer ownership table. A row
for business day 10 may arrive on day 20, after the day-15 decision. A
business-date filter then exposes a revision unavailable to the decision
maker. Joint ownership creates another ambiguity: copying an account
balance to both owners changes the economic unit and can double exposure.
These are observation-model questions before any causal estimator is called.

Five source-owner fixtures are specified separately from the interpretation
code. A current trusted arrival-clock statement excludes the late revision
and gives balance 90. Conflicting trusted arrival/business-clock statements
retain balances 90 and 110; missing clock documentation does the same. A
customer request with two joint owners and no allocation rule abstains.
Finally, a plausible but incorrect accepted migration specification yields
the unique wrong balance 110. The algorithm reports conditional acceptance
in that case, and the external fixture evaluation records wrong acceptance.
The result is not hidden by defining every supplied answer as correct.

Verified provenance means that the stated owner supplied the assertion, not
that the assertion's economic meaning is true. Schema profiles can generate
candidate clocks and entity relations; they cannot establish first
availability or ownership allocation by themselves. An analyst question
must therefore request the missing semantic fact and identify its source.
If pilot outcomes rank useful questions, the ranking and answers are frozen
before final evaluation. The executed final linkage experiment below uses
predeclared strata and independent audit information, with no final-outcome
question selection.

\subsection{Whole candidate maps and a recoverable mixture}
A scalar replacement rate is inadequate when plausible candidate entities
have different responses. Use two candidate portfolio histories at each
origin, with the entire response vector assigned consistently under each
map. In linkage stratum $s$, the two histories are correctly ordered with
probability $p_s$ and swapped otherwise. This is a one-to-one pair assignment,
not independent per-column linking. The match event is independent of the
complete potential histories conditional on the declared stratum and
matching information. That strong noninformative-linkage restriction is
part of identification; a name-match score is not evidence for it.
The executed map is redrawn by origin, while retaining all horizons within
that origin together. It models an independently reconstructed candidate
window, not a persistent customer alias across the entire calendar. A fixed
erroneous crosswalk or time-dependent matching mechanism would need its own
joint map law and covariance; the following evidence does not qualify them.

Let $\theta_{sA},\theta_{sB}$ be the true response vectors of the two
portfolios to the same price innovation. The two observable candidate
response vectors satisfy
\begin{equation}
 \begin{pmatrix}\eta_{s1}\\\eta_{s2}\end{pmatrix}
 =\begin{pmatrix}p_s&1-p_s\\1-p_s&p_s\end{pmatrix}
   \begin{pmatrix}\theta_{sA}\\\theta_{sB}\end{pmatrix},\qquad
 \theta_{sA}=\frac{p_s\eta_{s1}-(1-p_s)\eta_{s2}}{2p_s-1}.
 \label{s6:pair-mixing}
\end{equation}
This corrects two different candidate responses rather than assuming every
wrong entity has zero effect. An independently audited set of pair matches
estimates each $p_s$. At $p_s=1/2$ the matrix has rank one: these mean
equations identify only the sum of the two responses. In the symmetric
pair-swap observation law, exchanging portfolio labels preserves the
observed law; more unlabelled histories cannot resolve that named-portfolio
ambiguity. Higher mixture moments under additional restrictions could still
identify more than these mean equations do. If probabilities are not
calibrated, plausible maps produce a set of targets, not a probability-
weighted point estimate justified by convenience.

The funding experiment fits candidate conditional means separately by
stratum using an independent pilot and the prespecified driver basis.
The potentially misassigned lagged balance is excluded from this basis.
Under the stationary Gaussian driver law, conditional means of every past
driver and its centered square given the current driver are linear and
quadratic, respectively. Summing the stable balance recursion therefore
keeps the candidate outcome means in the driver basis after that lagged
balance is marginalized. This justification depends on the stated law;
it does not license dropping a confounding history in a new bank dataset.
Treatment residuals, candidate outcomes and stratum indicators produce
the candidate response moments. The conditional mean of treatment does
not depend on stratum, but allowing stratum interactions in its learner
preserves the common nuisance interface. Estimating candidate outcomes with
one pooled stratum mean could destroy the orthogonality needed for
stratum-specific inference; the executed design retains those interactions.

\subsection{Uncertainty from dates and calibration}
For fixed target weights $w_s$, let
$\theta=\sum_sw_s\theta_{sA}$. Stack all candidate response influences by
calendar date; the same dates, horizons and candidate histories induce
cross-component covariance $V_\eta$. With independent audits of fixed sizes
$M_s$, delta differentiation gives
\begin{align}
 d\theta_{sA}
 &=\frac{p_s}{2p_s-1}\,d\eta_{s1}
   -\frac{1-p_s}{2p_s-1}\,d\eta_{s2}
   +\frac{\eta_{s2}-\eta_{s1}}{(2p_s-1)^2}\,dp_s,\nonumber\\
 V_\theta
 &=J_\eta V_\eta J_\eta'
   +\sum_s j_sj_s'\frac{p_s(1-p_s)}{M_s},\qquad
 j_s=w_s\frac{\eta_{s2}-\eta_{s1}}{(2p_s-1)^2}.
 \label{s6:linkage-variance}
\end{align}
The calibration term is a matrix outer product: a single map probability
affects the entire response curve. Both $T$ and $M_s$ must grow in the
regular joint-sample approximation; a fixed small audit is not washed out
by arbitrarily many production records. Bias in the audit changes the
inverse map itself. Covariance correction cannot remove that bias.
If the Wilson audit interval contains $1/2$, the regular reference withholds the
regular inverse-map interval. This is an explicit availability rule, not
a claim of uniform post-selection coverage.
Section~\ref{s61:linkage-sets} supplies a separate projection procedure that
retains unbounded and disconnected financial confidence sets near rank
failure. It does not extend this delta approximation by changing its label.

In the registered experiment $p_0=.95$, $p_1=.8$, the true second-portfolio
price coefficient is .5 rather than the focal portfolio's 2, and both
histories share common shocks. With 720 evaluation dates and 100 audited
pairs per stratum, all 400 repetitions are available and financial coverage
is .9625; at 400 audits it is .950. Treating calibration as fixed gives
.960 and .950 respectively. Thus calibration is propagated, but its measured
variance contribution is modest relative to the calendar noise in this
design. The naive first-candidate analysis has bias about $-68{,}000$ and
$-80{,}000$ currency units in the two runs, with coverage .900 and .9075.
The correctly mapped financial target is about $375{,}149$ under the fixed
accounting inputs. These wide intervals are evidence about a difficult
dynamic observation problem, not a promised precise business case.

This construction is a bounded project derivation. Audit-based mismatch
correction provides one related approach \citep{slawski2025}; joint
probabilistic linkage and causal modeling offers a different route
\citep{guha2022} when a
well-specified joint likelihood and linking covariates are defensible. Such
a joint model may let outcomes inform linkage. Its posterior target and
feedback assumptions differ from this independent-audit frequentist
experiment; posterior match draws cannot simply be treated as additional
independent customers. The product records the selected mapping law, audit
population, covariance and unresolved semantic alternatives together.


The publisher version of \citet[Section 2.3, equation (2.4)]{guha2022}
models a constrained one-to-one map $z$, agreement fields and outcome
regressions jointly. The linked-record likelihood includes
$f_1(y_i\mid x_j,a_j)$; unlinked outcomes have $f_2(y_i)$. Thus the
posterior odds of a map contain the outcome likelihood ratio as well as
field agreement. Section 2.1 requires strong treatment ignorability and
positivity; the model cannot repair hidden action confounding by better
linkage. The journal version explicitly separates linking variables from
causal covariates and discusses conditioning/altering the likelihood when
they overlap. Its linked-population treatment effect and within-file
no-duplicates assumption differ from a repeated banking history with
joint ownership. These differences are why a production adapter must
represent the map law and target population, rather than passing arbitrary
posterior links into an otherwise unchanged temporal estimator.
